\section{Verification purpose}

\textbf{Case ID:} \texttt{\CaseID}. This case verifies the observable coupling of topographic slope into directional surface-fire spread through the standard ELMFIRE executable. It is the portable successor to Adam Laird's no-wind-with-slope component scenario and retains a flat control.

\section{Mathematical formulation and numerical implementation}

ELMFIRE combines the no-wind spread rate with a slope factor before supplying a direction-dependent normal speed $V_n(\theta)$ to $\partial_t\phi+V_n|\nabla\phi|=0$. With zero wind and zero slope, rotational symmetry is expected. A uniform nonzero slope breaks that symmetry along the fall-line axis.

\section{Derivation of expected behavior}

The flat control has equal forcing in all azimuths and therefore equal horizontal and vertical extents. The \SlopeAngle{} variant changes only slope and aspect; it must differ from the control and its final footprint must be elongated. These are metamorphic consequences of the governing coupling and do not require duplication of ELMFIRE's internal slope formula.

\section{Verification metrics and acceptance criteria}

Flat-control extent anisotropy must not exceed 0.05. The slope footprint's major/minor extent ratio must be at least 1.10. Normalized mean absolute arrival-time difference between slope and flat variants must be at least 0.05. The first metric detects grid bias, while the latter two show that slope is both active and directional.

\section{Simulation configuration and predicted outcome}

Both variants use a 2000 m square domain at \CellSize{}, Fuel Model 10, 3/5/6 percent dead-fuel moisture, zero wind, and a central point ignition. The slope variant uses 30 degrees and aspect 270 degrees; the control uses zero slope. The slope result is predicted to be anisotropic and measurably different from the circular control.

\section{Actual simulation results}

The generated status is \textbf{\VerificationStatus}; \PassingMetricCount{} of \MetricCount{} metrics pass across \VariantCount{} variants. Missing outputs produce NOT EVALUABLE.
\begin{figure}[H]\centering\IfFileExists{../figures/verification_summary.pdf}{\includegraphics[width=\textwidth]{verification_summary.pdf}}{\fbox{No result figure is available.}}\caption{Flat-control, slope, and difference arrival-time fields.}\end{figure}

\section{Calculated metrics and verification decision}

\begin{center}\begin{tabular}{@{}llll@{}}\toprule Metric & Expected & Calculated & Status \\\midrule \MetricRows \bottomrule\end{tabular}\end{center}
The authoritative decision is \textbf{\VerificationStatus} from \texttt{outputs/metrics.json}.

\section{Technical interpretation}

An anisotropic control indicates grid-orientation bias. A near-unity slope footprint or negligible field difference indicates that slope/aspect inputs were not consumed or not coupled to spread. Unexpected rotation should prompt inspection of aspect convention and raster orientation.

\section{Scope and limitations}

This case verifies directional response, not the absolute value of the slope coefficient and not agreement with an experimental fire. Existing CASE22 and CASE24 exercise more complex terrain interactions; this case isolates the slope-only executable response.
